\section{Analytical bounds for complete-source evaluation}
\label{sec:gt-source-certification}

The preceding results are analytical statements about the complete field.
This appendix explains how a finite calculation can supply their local
hypotheses without replacing the source. The interpolation and contour
bounds are proved first. The recorded computation at the end is supporting
evidence for one fixed tube, not a premise of the general theorems.

\subsection{An exact xi convolution and a lattice error bound}

For \(t\ge0\), the double-exponential source decay and the Gaussian integral
give
\begin{equation}
 H_t(z)=\frac1{\sqrt\pi}\int_{\R}\e^{-r^2}
                     H_0(z+2\ii\sqrt t\,r)\,\dd r .
 \label{eq:gt-xi-convolution}
\end{equation}
Indeed, after inserting \eqref{eq:field}, the absolute integrand is bounded
by a constant times
\[
 \Phi(u)\e^{|\Im z||u|}\e^{-(r+\sqrt t\,u)^2+tu^2}.
\]
Its integral is finite. Fubini and completing the square therefore
multiply the Fourier source by \(\e^{tu^2}\), exactly as required.
The time-zero formula already proved in Section~\ref{sec:source} may use
either \(\xi((1-\ii z)/2)/8\) or \(\xi((1+\ii z)/2)/8\), since
\(\xi(s)=\xi(1-s)\).

For the local application set
\[
 a=\frac{140101}{10},\qquad g=\e^{\pi a/8},\qquad
 \mathcal H_t=gH_t,\qquad 0\le t\le\frac1{10}.
\]
The constant gauge changes neither the pair coordinates nor the drift.
It cancels the principal gamma decay at this height. Write
\(f_{t,z}(r)=g\e^{-r^2}H_0(z+2\ii\sqrt t\,r)\).

\begin{lemma}[Source majorants]
\label{lem:gt-xi-majorants}
Let \(\alpha=9999/10000\), \(\beta=131/100\), and \(C=10^{13}\).
For \(|z-a|\le2\), \(r=\nu+\ii\eta\), \(|\eta|\le2\),
\begin{equation}
 |f_{t,z}(\nu+\ii\eta)|
 \le C\e^{\eta^2-\alpha\nu^2+\beta|\nu|}.
 \label{eq:gt-strip-majorant}
\end{equation}
On the larger spatial disk,
\begin{equation}
 |\mathcal H_t(a+w)|\le
 M_5:=\frac{2\cdot10^{16}}{\sqrt\alpha}
             \e^{\beta^2/(4\alpha)},\qquad |w|\le5 .
 \label{eq:gt-radius-five}
\end{equation}
\end{lemma}
\begin{proof}
Put \(q=\sqrt t\), \(s_0=(1+\ii(z+2\ii qr))/2\), and reflect \(s_0\)
to \(s\) with \(\Re s\ge1/2\). For \(\omega=s/2=X+\ii Y\) and
\(U=|\nu|\), the smaller disk and strip give
\[
 \frac14\le X\le\frac34+\frac4{25}U,\quad
 \left||Y|-\frac a4\right|\le\frac12+q,\quad
 3501<|Y|<3504,
\]
and hence \(|\omega|\le3505+4U/25\),
\(|s-1|\le|s|\le7010+8U/25\).
The right-half-plane Binet formula \cite[Eq.~5.9.10\_2]{dlmf} has remainder
\[
 R(\omega)=\int_0^\infty \e^{-\omega v}
       \left(\frac1{\e^v-1}-\frac1v+\frac12\right)\frac{\dd v}{v}.
\]
The parenthesis lies between zero and \(v/12\). To see both bounds,
put \(b=v/2\) and expand \(b\cosh b-\sinh b\) and
\((1+b^2/3)\sinh b-b\cosh b\): their coefficients at
\(b^{2n+1}\) are respectively \(2n/(2n+1)!\) and
\(2n(2n-2)/(3(2n+1)!)\), with the initial zero terms understood.
Thus \(|R(\omega)|\le1/(12X)\). Taking real parts and using
\(|Y|\arctan(X/|Y|)\le X\) yields
\[
 |\Gamma(\omega)|\e^{\pi|Y|/2}
 \le\sqrt{2\pi}\,|\omega|^{X-1/2}\e^{1/(12X)}.
\]
Euler summation gives
\(\zeta(s)=s/(s-1)-s\int_1^\infty\{v\}v^{-s-1}\,\dd v\);
the present \(\Re s\ge1/2\), \(|\Im s|>7002\) imply
\(|\zeta(s)|\le3|s|\). Use the completion
\(\xi(s)=s(s-1)\pi^{-s/2}\Gamma(s/2)\zeta(s)/2\).
Since \(g\e^{-\pi|Y|/2}<4\), \(\sqrt{2\pi}<3\) and
\(\e^{1/3}<3/2\), the integrand is at most
\[
 \frac{4\cdot3\cdot(3/2)\cdot3}{16}
 (7010+8U/25)^3(3505+4U/25)^{1/4+4U/25}\e^{\eta^2-U^2}.
\]
Its prefactor at zero is below \(10^{13}\).
Expanding the logarithm with \(\log(1+x)\le x\), the extra linear
coefficient is less than \(131/100\), and the extra quadratic
coefficient is less than \(1/10000\).
These comparisons follow from
\(\log3505<409/50\), for which
\(\sum_{j=0}^{24}(409/50)^j/j!>3505\).
This proves \eqref{eq:gt-strip-majorant}.

For \(|w|\le5\) and real \(r\), replace these bounds by
\(X\le3/2+4|r|/25\), \(|\omega|\le3506+4|r|/25\),
\(|s|\le7012+8|r|/25\), and
\(g\e^{-\pi|Y|/2}<8\).
The same calculation gives
\(10^{16}\e^{-\alpha r^2+\beta|r|}\).
Integration in \eqref{eq:gt-xi-convolution}, followed by completion
of the square on each half-line, proves \eqref{eq:gt-radius-five}.
\end{proof}


\begin{corollary}[Majorants at moving boundary centers]
\label{gt-boundary-majorants}
For every \(b\in[13840,14180]\), put \(g_b=\e^{\pi b/8}\).
The bounds \eqref{eq:gt-strip-majorant} and \eqref{eq:gt-radius-five}
hold with \(a,g\) replaced by \(b,g_b\), with the same
\(C,\alpha,\beta,M_5\).
\end{corollary}
\begin{proof}
For the radius-two disk and \(|\Im r|\le2\), put \(V=|\Re r|\).
The coordinates in the preceding proof now satisfy
\[
 \frac14\le X\le\frac34+\frac4{25}V,\quad
 \left||Y|-\frac b4\right|\le\frac12+\sqrt t,\quad
 3459<|Y|<3547.
\]
Hence \(|\omega|\le3550+4V/25\) and
\(|s-1|\le|s|\le7100+8V/25\). The Binet and Euler estimates
still apply, and \(g_b\e^{-\pi|Y|/2}<4\).
The constant prefactor is at most \(27(7100)^3<10^{13}\);
indeed \(3550^{1/4}<8\). The extra log-linear coefficient is at most
\[
 \frac4{25}\frac{409}{50}
 +\frac{(1/4)(4/25)}{3550}
 +\frac{3(8/25)}{7100}<\frac{131}{100},
\]
and the quadratic coefficient is \((4/25)^2/3550<1/10000\).
The exact inequality
\(\sum_{j=0}^{24}(409/50)^j/j!>3552\) proves the logarithm bound
used here and below. This proves the strip estimate.

For the radius-five disk and real \(r\), use
\(X\le3/2+4|r|/25\), \(3458<|Y|<3547\),
\(|\omega|\le3552+4|r|/25\), \(|s|\le7104+8|r|/25\)
and \(g_b\e^{-\pi|Y|/2}<8\).
The prefactor is at most
\((27/4)(7104)^3(3552)=8595785870475264<10^{16}\).
The log-linear coefficient is bounded by
\[
 \frac4{25}\frac{409}{50}
 +\frac{4/25}{3552}+\frac{3(8/25)}{7104}<\frac{131}{100},
\]
and the quadratic coefficient by \((4/25)^2/3552<1/10000\).
Integration and completion of the square give the same \(M_5\).
\end{proof}

For an affine boundary \(b(t)=B+v(t-1/20)\) in this interval of
centers, the fixed gauge \(g_B\) and local gauge obey
\[
 \frac{g_B}{g_{b(t)}}\le A_v:=\e^{\pi|v|/160}
 \qquad(0\le t\le1/10).
\]
Consequently the full affine derivative formula \eqref{gt-moving-boundary-derivative}
may use \(M_*=A_vM_5\) and \(R_*=5\).
The lattice allowance below likewise holds at every such center with
local gauge \(g_b\); multiplication by \(A_v\) pays a fixed gauge.
The gauge remains constant when differentiating in time.

\begin{proposition}[Finite lattice payment]
\label{prop:gt-lattice-payment}
For \(h>0\) with \(12/h\) integral, let
\[
 \mathcal H_t^{[h]}(z)=\frac h{\sqrt\pi}
                    \sum_{|kh|\le12}f_{t,z}(kh).
\]
On \(|z-a|\le2\) throughout the stated time interval,
\begin{equation}
 |\mathcal H_t-\mathcal H_t^{[h]}|
 \le
 \frac{4C\e^{4+\beta^2/(4\alpha)}}
      {\sqrt\alpha(\e^{4\pi/h}-1)}
 +\frac{2C\e^{-144\alpha+12\beta}}
      {\sqrt\pi(24\alpha-\beta)}=:E_h .
 \label{eq:gt-lattice-error}
\end{equation}
Finite-sum arithmetic errors must be added separately.
\end{proposition}
\begin{proof}
For \(\widehat f(\omega)=\int f(r)\e^{-\ii\omega r}\,\dd r\),
shift the line to imaginary height \(-2\) when \(\omega>0\) and to
\(2\) when \(\omega<0\). The vertical ends vanish by
\eqref{eq:gt-strip-majorant}, and
\[
 |\widehat f(\omega)|\le
 2C\sqrt{\pi/\alpha}\,\e^{4+\beta^2/(4\alpha)}\e^{-2|\omega|}.
\]
The Fourier coefficients of the uniformly convergent periodization
\(h\sum_{k\in\mathbb Z}f(x+kh)\) are
\(\widehat f(2\pi n/h)\). They are absolutely summable; evaluating at
zero and summing their geometric tails gives the first term of
\eqref{eq:gt-lattice-error}. On \([12,\infty)\) the real majorant
decreases, so the omitted lattice values, times \(h\), are bounded by
twice its tail integral. The concave exponent lies below its tangent
at \(12\), giving the second term.
\end{proof}

\subsection{One polynomial and two independently paid radii}

The interpolation argument is useful in a general form. Suppose a true
heat field \(\mathcal H\) obeys
\(|\mathcal H_t(a+w)|\le M_S\) for \(|w|\le S\), \(t\in[u,v]\), where \(u<v\).
Choose \(0<\rho<S\), \(0<R<S\), integers \(0\le d<M\), \(n\ge1\), and sample at
\(w_j=\rho\e^{2\pi\ii j/M}\) and Lobatto times
\(t_l=(u+v)/2+\Delta\cos(\pi l/n)\), \(\Delta=(v-u)/2\).
Assume every sample error is at most \(\varepsilon_s\).
Take its discrete Fourier coefficients through degree \(d\), then
interpolate those spatial polynomials at the \(n+1\) time nodes.
Call the resulting exact polynomial \(P(t,w)\).

\begin{proposition}[Uniform tensor error]
\label{prop:gt-tensor-error}
For \(|w|\le R\), \(t\in[u,v]\), one has
\(|\mathcal H_t(a+w)-P(t,w)|\le E(R)\), where
\begin{align}
 E(R)&:=(2n+1)\left[
 \varepsilon_s\sum_{k=0}^d(R/\rho)^k
 +M_S\frac{(\rho/S)^M}{1-(\rho/S)^M}
                \sum_{k=0}^d(R/S)^k
 +M_S\frac{(R/S)^{d+1}}{1-R/S}\right]\notag\\
 &\quad+
 \frac{M_S(2n+2)!}{(n+1)!(S-R)^{2n+2}}\,
            2^{1-n}\Delta^{n+1}.
 \label{eq:gt-tensor-error}
\end{align}
Coefficient rounding, if treated as exact values rather than enclosed
balls, needs its own additive uniform payment.
\end{proposition}
\begin{proof}
At a time node, write
\(\mathcal H(a+w)=\sum_{k\ge0}a_kw^k\).
Cauchy gives \(|a_k|\le M_S/S^k\).
The root-of-unity sum for coefficient \(k<M\) equals
\(\sum_{j\ge0}a_{k+jM}\rho^{k+jM}\).
Consequently sample propagation, aliasing and discarded spatial powers
give respectively the three bracketed terms.
For Lobatto interpolation the discrete cosine coefficients have modulus
at most twice the largest node value. The endpoint-half-weight cosine
expansion and \(|T_k(x)|\le1\) give operator norm at most \(2n\);
we use \(2n+1\).

The remaining error interpolates the true field in time.
Its heat equation implies
\(\partial_t^{n+1}\mathcal H=(-1)^{n+1}\partial_z^{2n+2}\mathcal H\).
Cauchy on radius \(S-R\) bounds this derivative by
\(M_S(2n+2)!/(S-R)^{2n+2}\).
The monic node product is
\(2^{1-n}(x^2-1)U_{n-1}(x)\), of modulus at most \(2^{1-n}\)
on \([-1,1]\). Scaling contributes \(\Delta^{n+1}\).
The real interpolation remainder divided by \((n+1)!\) proves the
last term; for complex values apply it to the real projection in the
direction of the error. No heat equation is imposed on \(P\).
\end{proof}

The same polynomial can be used at several radii, provided each radius
has its own \(E(R)\). With \(\rho=9/10\), \(S=5\), \(M=128\) and
\(d=n=32\), this gives both the original collar radius \(21/20\)
and the enlarged exclusion radius \(8/5\). The larger error belongs
to the larger disk; the original smaller error remains available for
the moment computation. This separation prevents a geometry extension
from needlessly degrading the coordinate estimates.

\subsection{Count, logarithm, coordinates and curvature}

Let \(r_-<\rho<r_+\). Suppose one sign \(\sigma\in\{-1,1\}\)
gives, on both entire boundary circles and every time,
\begin{equation}
 \Re\frac{\sigma\mathcal H_t(a+w)}{w^2}\ge m_\pm>0,\qquad
 \left|\frac{\sigma\mathcal H_t(a+w)}{w^2}\right|\le M_\pm .
 \label{eq:gt-half-plane}
\end{equation}
Polynomial interval bounds pay this condition after subtracting
\(E(r_\pm)/r_\pm^2\) from a lower real-part bound and adding it to a
modulus upper bound. Bounds over closed time cells and angular arcs,
not just their midpoints, are required.

The normalized boundary loop lies in the right half-plane and has
winding zero, so \(\mathcal H_t\) has exactly two zeros inside each
circle. Their equal counts exclude every zero in the annulus.
The harmonic minimum principle preserves the positive real part there,
giving a principal logarithm on an open collar. For a fixed real
centering \(\ell_*\), set
\[
 \mathcal L_{\mathrm{ann}}(t,w)=
 \operatorname{Log}\bigl(\sigma\mathcal H_t(a+w)/w^2\bigr)-\ell_*,
 \quad
 B_\pm^2=
 \max\{|\log m_\pm-\ell_*|,|\log M_\pm-\ell_*|\}^2+\pi^2/4.
\]
Then \(|\mathcal L_{\mathrm{ann}}|\le B_\pm\) on the corresponding boundaries.
Here \(\mathcal L_{\mathrm{ann}}\) denotes this local logarithm, not a Laguerre expression.

For completeness, if \(V\) is analytic on a neighborhood of the annulus,
with boundary bounds \(A_\pm\), its \(M\)-node angular mean on \(\rho\)
differs from its exact mean by at most
\begin{equation}
 T_M(A_-,A_+)=
 A_-\frac{(r_-/\rho)^M}{1-(r_-/\rho)^M}
 +A_+\frac{(\rho/r_+)^M}{1-(\rho/r_+)^M}.
 \label{eq:gt-annular-error}
\end{equation}
Indeed, expand \(V\) in Laurent powers. Cauchy controls positive powers
from the outer circle and negative powers from the inner circle.
Only multiples of \(M\) survive the discrete average, and summing them
gives \eqref{eq:gt-annular-error}.

Write \(w=z-a\), \(\gamma=c-a\), \(Q=(w-\gamma)^2-D\), and let brackets
denote the normalized angular mean. Contour integration by parts gives
\begin{equation}
 p_j=-j\langle w^j\mathcal L_{\mathrm{ann}}\rangle,\quad
 \gamma=p_1/2,\quad D=p_2/2-\gamma^2,\quad
 K=\left\langle\frac{2w(w-\gamma)\mathcal L_{\mathrm{ann}}}{Q^2}\right\rangle.
 \label{eq:gt-log-moments}
\end{equation}
For the last equality, the logarithmic derivative of the pair contributes
\(\oint Q'/Q^2=0\); the term \(2/(wQ)\) also has zero contour integral,
since all its poles are inside and it decays cubically at infinity.
Integrate the remaining \(\mathcal L_{\mathrm{ann}}'/Q\) by parts. All exact means in
\eqref{eq:gt-log-moments} annihilate a constant log centering.

Apply \eqref{eq:gt-annular-error} to the moment integrands with
\(A_\pm=j r_\pm^jB_\pm\). Enclose both moments on each time cell and
propagate their intervals into \(\gamma,D\). If \(C_*\ge|\gamma|\),
\(D_*\ge|D|\), a sufficient denominator bound on radius \(r>C_*\) is
\[
 |Q|\ge q_r=(r-C_*)^2-D_*>0.
\]
A sharper direct bound can replace it. The drift integrand then has
boundary bounds
\(A_\pm=2r_\pm(r_\pm+C_*)B_\pm/q_{r_\pm}^2\).
Thus the actual coordinate intervals and the annular tail together
enclose the drift, including when the gap interval contains zero.

An approximate field with contour error \(\delta<m\), where \(m\)
is its boundary modulus floor, also gives the aligned log error
\(\ell=-\log(1-\delta/m)\). This follows from the convergent logarithm
series on a collar. Adding angular errors \(q_1,q_2\) gives
\[
 |c-\widetilde c|\le e=\frac{\rho\ell+q_1}{2},\qquad
 |D-\widetilde D|\le d=\rho^2\ell+q_2/2+2(\rho-b)e
\]
when both pairs have buffer \(b\). If only the true center is bounded,
replace the last factor by \(2(\rho-b)+e\).
For the drift estimate retain both pair buffers \(b\). With
\(J_*=(4\rho-2b)e+d\), a centered approximate log bounded by
\(\Lambda_*\) yields
\begin{equation}
 |K-\widetilde K|\le
 \frac{2\rho\ell}{b^3}
 +\rho\Lambda_*\left(\frac{2e}{b^4}
       +\frac{4(2\rho-b)J_*}{b^6}\right)+q_K .
 \label{eq:gt-value-transfer}
\end{equation}
To verify the kernel term, use \(|Q|,|\widetilde Q|\ge b^2\),
\(|Q-\widetilde Q|\le J_*\) and
\[
 |(1/Q)'-(1/\widetilde Q)'|
 \le2e/b^4+4(2\rho-b)J_*/b^6.
\]
The log-error term uses \(|(1/Q)'|\le2/b^3\).
This proves \eqref{eq:gt-value-transfer} with no inverse internal gap.

\subsection{Supporting evidence on one fixed tube}

The recorded interval-arithmetic application uses the same complete-source
array, the analytical bounds above and outward rounding
\cite{gtcertificate,gtgeometry}. It gives the following enclosures;
the distinct roles of terminal and whole-interval coordinates matter.
\begin{center}
\small
\begin{tabular}{@{}p{0.40\textwidth}p{0.51\textwidth}@{}}
\toprule
Quantity & Recorded conclusion on \(0\le t\le1/10\)\\
\midrule
Original geometry & Exactly two zeros inside \(3/4\), zero-free collar to \(21/20\).\\
Full drift & \(-0.615518848\le K(t)\le-0.599329616\).\\
Incoming gap & \(0.17873684703288\le D(1/10)<0.17873685811\).\\
Signed passage margin using \(K\le-0.59\) & Greater than \(0.00058487578\).\\
Direct endpoint & \(D(0)\in[0.00142117117778,\,0.00142118228384]\).\\
Enlarged geometry & Exactly the same two zeros inside radius \(8/5\).\\
New whole-interval coordinates & \(|c-a|<0.063763113\), \(|D|<0.179116620\).\\
\bottomrule
\end{tabular}
\end{center}
The enlarged geometry uses the same cached samples and evaluates no
new drift. It instantiates the grouped-divisor sign theorem independently
of the old measured drift band. More explicitly, at
\(a=140101/10\), \(C=7/100\) and \(d=1/5\), the auxiliary radius
\(8/5\) gives
\[
 X=\frac{\sqrt{39}}5-\frac7{100}>\frac{113}{100},
 \qquad X^2>1+d.
\]
The same two occurrences also lie inside radius \(9/10\), since the
original collar is zero-free. The older radius \(21/20\) gives only
\(X=\sqrt{41}/20-7/100<1\), so it does not pay the sufficient
disk-to-horizontal test. The reflected target pair contributes the
lower allowance \(\kappa_0=5000/1962848636289\). At \(v=1/10\), its
constant-interaction transport allowance satisfies
\[
 2\int_0^v \e^{-4\kappa_0(v-t)}\,\dd t
 \ge \frac15-\frac{\kappa_0}{25}>\frac{179}{1000},
\]
by \(\e^{-s}\ge1-s\). This exceeds the recorded terminal gap.
The reflected group therefore proves a sign but does not by itself
pay the useful cost threshold; other exterior mass can reduce the
actual cost.

A later supporting acquisition implements the real-bracket mechanism
of Section~\ref{gt-source-curvature} \cite{gtbrackets}. All 176 fixed
moving real brackets have opposite, nonzero boundary signs throughout
\(0\le t\le1/10\). The boundary signs are enclosed using degree-12
time interpolation with the full affine chain-rule remainder, complete
source lattice and tails, and interval arithmetic. Same-time disjointness
supplies at least one distinct real exterior occurrence per bracket on
each slice; no singleton count, simplicity or persistent label is claimed.
Independently replayed target coordinates on 512 time cells give
\[
 M(t)>0.5937558081007126288>0.59\qquad(0\le t\le1/10).
\]
Only the reflected target pair is separately retained; all other omitted
conjugate groups are favorable. The same brackets with globally collapsed
coordinate information gave lower allowances approximately
\(0.559912322830\) and \(0.582325753776\), which did not reach the
threshold. Those failed allowances are not upper bounds on actual mass.
The successful certificate preserves time-local coordinate information.
Its positive lower bound on \(D\) was already supplied independently,
so the new result acquires exterior magnitude rather than new target
reality. The earlier measured drift band is stronger. These are recorded
supporting calculations, internally assessed with one cache replay; the
present manuscript integration performs no new source evaluations.

The absolute-value cost is at least \(2/10=1/5\), exceeding the incoming
gap; the signed cost succeeds. More explicitly, \(K\le-0.59\)
gives \(W(0)\le\e^{-0.236}<0.789780673932803\) in
\eqref{eq:gt-backward-gap}, hence \(D(0)>0.00074055469\).
The direct endpoint enclosure is stronger than this transported bound. Moreover \(K\in[-5/8,0]\) and the
inner radius \(3/4\) give \(D'\ge19/32\), so the direct endpoint also
encloses the interval minimum. All these numerical conclusions concern
one fixed tube and retain their internal-review status. They supply no
height-uniform source inequality, new global zero coverage or exclusive
transport advantage.

As an exact independent control, \(F_t(a+w)=w^2-2(t-1/20)\) passes
the same count and logarithm conditions but has \(K=0\).
Its terminal gap is \(1/10\), its full-interval cost is \(1/5\),
and its strict margin is \(-1/10\). Thus the supporting procedure can
accept geometry at a double while correctly rejecting passage.
